\usepackage[T1]{fontenc}
\usepackage[utf8]{inputenc}
\usepackage{lmodern}
\usepackage{amsmath,amssymb,amsthm,mathtools}
\usepackage{esint}      % supplies \fint
\usepackage[margin=1.05in]{geometry}
\usepackage{enumitem}
\usepackage{booktabs}
\usepackage{longtable}
\usepackage{microtype}
\usepackage{seqsplit}
\usepackage{xcolor}
\usepackage[colorlinks=true,linkcolor=blue,citecolor=blue,urlcolor=blue]{hyperref}

% Long Lean identifiers may break across lines without inserting a hyphen.
\newcommand{\code}[1]{\texttt{\seqsplit{#1}}}
\setlength{\emergencystretch}{3em}

% ----------------------------------------------------------------------------
%  Blueprint macros.
%
%  These are defined here as no-ops so that this file compiles with plain
%  pdflatex.  When this document is imported into a `leanblueprint` project,
%  DELETE this block: plasTeX supplies the real definitions.
%
%    \uses{lab1,lab2}  dependency edges in the blueprint DAG
%    \lean{Decl.Name}  the Lean declaration realising the item
%    \leanok           the item is fully formalised
%    \notready         the statement is not yet in a formalisable shape
%    \discussion{n}    link to a GitHub issue
% ----------------------------------------------------------------------------
\providecommand{\uses}[1]{}
\providecommand{\lean}[1]{}
\providecommand{\leanok}{}
\providecommand{\notready}{}
\providecommand{\discussion}[1]{}

% ----------------------------------------------------------------------------
%  Theorem environments.  One shared counter, so that every citable item in
%  the document has a unique number and there is no ambiguity about what
%  "Lemma 7.3" means.
% ----------------------------------------------------------------------------
\theoremstyle{plain}
\newtheorem{theorem}{Theorem}[chapter]
\newtheorem{proposition}[theorem]{Proposition}
\newtheorem{lemma}[theorem]{Lemma}
\newtheorem{corollary}[theorem]{Corollary}

\theoremstyle{definition}
\newtheorem{definition}[theorem]{Definition}
\newtheorem{convention}[theorem]{Convention}
\newtheorem{notation}[theorem]{Notation}
\newtheorem{constants}[theorem]{Constants}

\theoremstyle{remark}
\newtheorem{remark}[theorem]{Remark}
\newtheorem{fnote}[theorem]{Formalisation note}
\newtheorem{test}[theorem]{Regression test}

% Intended Lean module for the current section.
\newcommand{\module}[1]{%
  \par\noindent{\small\raggedright\textsf{Lean module:}\ \texttt{#1}\par}\smallskip}

% ----------------------------------------------------------------------------
%  Notation
% ----------------------------------------------------------------------------
\newcommand{\R}{\mathbb R}
\newcommand{\N}{\mathbb N}
\newcommand{\Z}{\mathbb Z}
\newcommand{\Sph}{\mathbb S}
\newcommand{\Ec}{\mathcal E}
\newcommand{\Hs}{\mathcal H}
\newcommand{\Lloc}{L^1_{\mathrm{loc}}}
\newcommand{\ind}[1]{\mathbf 1_{#1}}
\newcommand{\Per}{P}
\newcommand{\rbd}{\partial^*}
\newcommand{\ebd}{\partial^{\mathrm e}}
\newcommand{\dvg}{\operatorname{div}}
\newcommand{\dvgt}{\operatorname{div}_\tau}
\newcommand{\Exc}{\operatorname{\mathbf E}}
\newcommand{\ExcB}{\operatorname{\mathbf E}^{\mathrm b}}
\newcommand{\Capa}{\operatorname{Cap}}
\newcommand{\cof}{\operatorname{cof}}
\newcommand{\Lip}{\operatorname{Lip}}
\newcommand{\dist}{\operatorname{dist}}
\newcommand{\spt}{\operatorname{spt}}
\newcommand{\Tr}{\operatorname{Tr}}
\newcommand{\avg}{\fint}
\DeclareMathOperator{\Id}{Id}
\newcommand{\halfsp}[1]{\{y : y\cdot #1 < 0\}}

\title{\bfseries The Three-Dimensional Liquid-Drop Threshold\\[2mm]
\large A Lean-friendly blueprint specification}
\author{Autoformalisation specification}
\date{}

